\section{Coding finite probability in iPRA}\label{sec:coding}

\subsection{Signed integers and rationals}

A signed integer is represented by a pair \((a,b)\in\N^2\), denoting
\(a-b\).  Equality and order are
\[
 (a,b)=(c,d) \quad\Longleftrightarrow\quad a+d=c+b,
\]
\[
 (a,b)\le(c,d) \quad\Longleftrightarrow\quad a+d\le c+b.
\]
Addition, multiplication, negation, and absolute value are primitive recursive.
For instance,
\((a,b)(c,d)=(ac+bd,ad+bc)\), and a decidable comparison selects the positive
or negative part for absolute value.

A rational code is \((z,q)\), where \(z\) is a signed-integer code and
\(q>0\).  It denotes \(z/q\).  Equality and order are again decided by positive
cross-multiplication.  Reduction to lowest terms is optional; no proof below
relies on unique representation.

\begin{lemma}[Common denominator normalization]\label{lem:common-den}
For every coded finite list of rational numbers
\(x_0,\ldots,x_{m-1}\), there is a primitive-recursively computed positive
integer \(D\) and signed integers \(a_i\) such that \(x_i=a_i/D\) for every
\(i<m\).
\end{lemma}

\begin{proof}
Take the product of the positive input denominators.  Multiplying each
numerator by the product of the remaining denominators gives the required
\(a_i\).  The bounded product and the transformed list are primitive recursive.
All displayed equalities follow by distributivity; no gcd computation is
needed.
\end{proof}

\subsection{Finite weighted laws}

\begin{definition}[Finite law]\label{def:finlaw}
A finite weighted law is coded by
\[
 m\in\N,\qquad w:\Fin(m)\to\N,
 \qquad W=\sum_{i<m}w_i>0.
\]
The probability of an event \(A\subseteq\Fin(m)\), represented by a decidable
indicator \(a_i\in\{0,1\}\), is
\[
 \Pp(A)=\frac{1}{W}\sum_{i<m}w_i a_i.
\]
A rational observable is a finite list \(X_i\in\Q\), and its expectation is
\[
 \E[X]=\frac{1}{W}\sum_{i<m}w_iX_i.
\]
\end{definition}

The use of natural weights loses no finite rational law.  If the listed
probabilities are \(p_i=r_i/s_i\), choose a common denominator \(W\), put
\(w_i=Wp_i\), and delete zero-weight padding if desired.  Conversely, natural
weights produce the rational probabilities \(w_i/W\).

\begin{lemma}[Primitive-recursive evaluator]\label{lem:pr-evaluator}
Given codes for a finite law and a rational observable, a code for
\(\E[X]\) is primitive recursive.  Equality and order between two such
expectations are quantifier-free after clearing the positive denominators.
\end{lemma}

\begin{proof}
Use \Cref{lem:common-den} to write \(X_i=a_i/D\).  Then
\[
 \E[X]=\frac{\sum_{i<m}w_i a_i}{WD}.
\]
The numerator is a bounded sum of primitive-recursive signed products, and
\(WD>0\).  Comparing two expectations is therefore a comparison of two signed
cross-products.
\end{proof}

\subsection{Events, null sets, and almost-sure equality}

\begin{definition}[Finite almost-sure equality]\label{def:ae}
For observables \(X,Y\) on the same finite law, write \(X=Y\) almost surely
when
\[
  \sum_{i<m} w_i\ind_{X_i\ne Y_i}=0.
\]
An event is null when its weighted indicator sum is zero.
\end{definition}

This definition is stronger computationally than merely saying that the
probability quotient equals zero: its numerator is explicitly zero.  Since all
summands are natural and nonnegative, zero-sum elimination gives
\(w_i\ind_{X_i\ne Y_i}=0\) for every \(i<m\).  ProveIt already contains exactly
this bounded-sum pattern: a zero bounded sum has zero terms, and conversely
pointwise zero terms give a zero sum.

\subsection{Compiling bounded algebra to quantifier-free induction}

\begin{proposition}[Finite-sum compilation pattern]\label{prop:sum-compile}
Each of the following has an explicit \(\ipra\) derivation by
quantifier-free induction on the bound:
\begin{align*}
 \sum_{i<n+1}f(i)&=\sum_{i<n}f(i)+f(n),\\
 \sum_{i<n}(f(i)+g(i))&=\sum_{i<n}f(i)+\sum_{i<n}g(i),\\
 \left[\forall i<n\;f(i)=g(i)\right]&\Longrightarrow
     \sum_{i<n}f(i)=\sum_{i<n}g(i),\\
 j<n&\Longrightarrow f(j)\le\sum_{i<n}f(i)
 \quad\text{when all values are natural.}
\end{align*}
The bounded universal premise is represented by a zero sum of decidable
failure indicators.
\end{proposition}

\begin{proof}
The first equation is the defining equation of the primitive-recursive sum.
For the second, the induction invariant is the displayed equality.  The base
case is \(0=0+0\); the successor step substitutes the first recurrence and
uses associativity and commutativity.  For congruence, maintain equality of the
partial sums and use the zero failure-count to obtain equality at the new
index.  The final inequality follows because the target term occurs as one
nonnegative summand; an induction on the upper bound gives the formal proof.
All invariants are quantifier-free equations or inequalities.
\end{proof}

\begin{corollary}[Finite Fubini and regrouping]\label{cor:finite-fubini}
For coded finite arrays,
\[
 \sum_{i<m}\sum_{j<n}a_{ij}
 =\sum_{j<n}\sum_{i<m}a_{ij}
 =\sum_{k<mn}a_{\lfloor k/n\rfloor,\,k\bmod n}.
\]
More generally, a finite sum may be regrouped by any primitive-recursive map
between finite index sets, provided multiplicity is tracked.
\end{corollary}

\begin{proof}
Induct on one rectangle dimension and use distributivity.  The mixed-radix
map \(k\mapsto(\lfloor k/n\rfloor,k\bmod n)\) is primitive recursive by division
with remainder.  The two directions of the bijection are verified by the
usual quotient-remainder equations, which are available in the arithmetic
library.
\end{proof}

\subsection{Rational inequalities without division}

Whenever a theorem contains fractions, the formal target is its
positive-denominator clearing.  For example, monotonicity
\[
 \frac{A}{WD}\le\frac{B}{WE}
\]
becomes \(AE\le BD\) after checking the correct denominators and signs.  A
proof should never invoke an abstract ordered field when a direct natural
identity is available.  This policy both shortens the object-language proof
and makes the computational content visible.
